%% presentation.tex
%% Beamer version of the umich-revealjs Quarto presentation
%%
%% Compile with:
%%   pdflatex presentation.tex   (two passes for cross-references)
%% or
%%   lualatex  presentation.tex  (if you want fontspec/Roboto)
%%
%% The beamerthemeumich.sty file must be in the same directory.

\documentclass[aspectratio=169, 12pt]{beamer}

\usetheme{umich}

% Short forms appear in the footline
\title{Title of your presentation}
\subtitle{Subtitle of your presentation}

\author[Name]{%
  Name \orcid{0000-0000-0000-0000} \\  % optional — remove \orcid to omit ORCID icon
  \texttt{\footnotesize johndoe@umich.edu}}


  \institute[University of Michigan]{%
Department of political science, University of Michigan}

\date{\today}

% --------------------------------------------------------------------------
% Extra packages for content
% --------------------------------------------------------------------------
\usepackage{listings}
\lstset{
  basicstyle=\ttfamily\footnotesize,
  backgroundcolor=\color{umichNavy!8!white},
  frame=single,
  framerule=0.4pt,
  rulecolor=\color{umichNavy!40!white},
  breaklines=true,
}

% --------------------------------------------------------------------------
\begin{document}

% ------ Title slide --------------------------------------------------------
\begin{frame}[plain]
	\titlepage
\end{frame}

% ------ Slide 1 ------------------------------------------------------------
\begin{frame}{Example slide}
	\framesubtitle{This is a subtitle}

	Here we have some text that may run over several lines of the slide frame,
	depending on how long it is.

	\begin{itemize}
		\item first item
		      \begin{itemize}
			      \item A sub item
		      \end{itemize}
	\end{itemize}

	\medskip
	Next, we'll briefly review some theme-specific components.

	\begin{itemize}
		\item Note that \textit{all} of the standard Beamer features can be used
		      with this theme, even if we don't highlight them here.
	\end{itemize}
\end{frame}

% ------ Slide 2 ------------------------------------------------------------
\begin{frame}{Additional theme classes}
	\framesubtitle{Some extra things you can do with the clean theme}

	Special classes for emphasis

	\begin{itemize}
		\item \alert{important note} — \texttt{\textbackslash alert\{\}} for
		      default emphasis (maroon)
		\item \textcolor{umichMaroon}{custom colour} — via
		      \texttt{\textbackslash textcolor\{umichMaroon\}\{\}}
		\item \colorbox{umichMaroon!20}{\textcolor{umichMaroon}{highlighted}} — via
		      \texttt{\textbackslash colorbox}
	\end{itemize}

	\medskip
	Cross-references

	\begin{itemize}
		\item Use \texttt{\textbackslash hyperlink\{label\}\{text\}} for internal
		      Beamer links — analogous to the \texttt{.button} class
	\end{itemize}
\end{frame}

% ------ Slide 4 ------------------------------------------------------------

\end{document}
